\chapter{Hamming--Wasserstein 距離への還元}
\label{ch:classical-reduction}

量子 $W_1$ が古典 Wasserstein 距離の拡張であるためには，計算基底で対角な状態に制限したとき，
元の距離を正確に再現しなければならない．本章ではこれを示す．

\section{離散 cube 上の古典 \texorpdfstring{$W_1$}{W1}}

\begin{definition}{Hamming 距離}{q-hamming}
  $x,y\in\{0,1\}^n$ に対し

  \[
    d_H(x,y):=\sum_{i=1}^n\mathbf{1}_{\{x_i\neq y_i\}}
  \]

  と定める．これは $x$ と $y$ が異なる座標の個数である．
\end{definition}

確率分布 $p,q$ の coupling 全体を $\Pi(p,q)$ とすると，Hamming 距離に対する古典 $W_1$ は

\[
  \Wass_1^H(p,q)
  :=\min_{\gamma\in\Pi(p,q)}
  \sum_{x,y\in\{0,1\}^n}d_H(x,y)\gamma(x,y)
\]

である．有限集合上の Kantorovich 双対性により，これは全質量が $0$ の符号付きベクトル上の
Kantorovich--Rubinstein ノルム

\[
  \norm{r}_{\rm KR}
  :=\max\left\{
    \sum_xf(x)r(x)
    \;\middle|\;
    \abs{f(x)-f(y)}\leq d_H(x,y)
  \right\}
\]

が誘導する距離であり，$\Wass_1^H(p,q)=\norm{p-q}_{\rm KR}$ である．
この有限次元線形計画の双対性は標準結果として用いる．

\begin{lemma}{一座標だけ異なり得る分布}{q-classical-neighbor}
  $p,q$ が第 $i$ 座標を除いた周辺分布について一致するならば

  \[
    \Wass_1^H(p,q)\leq1
  \]

  である．
\end{lemma}

\begin{proof}
  残りの $n-1$ 座標を $z$ と書き，共通の周辺分布を $m(z)$ とする．
  まず $z$ を $m$ に従って同じ値に coupling し，その条件下で第 $i$ 座標の
  条件付き分布どうしを任意に coupling する．得られる二点は第 $i$ 座標以外で一致するので，
  Hamming コストは常に $1$ 以下である．
\end{proof}

\section{計算基底の点}

量子 $W_1$ の局所性から，まず Dirac 分布に対応する純粋状態の距離を計算する．

\begin{lemma}{一 qubit 周辺状態による下界}{q-marginal-lower}
  $X\in\mathcal{O}_n^0$ とし，$X_i:=\Tr_{\{1,\ldots,n\}\setminus\{i\}}X$ を
  第 $i$ qubit 上の縮約とする．このとき

  \[
    \norm{X}_{W_1}\geq\frac12\sum_{i=1}^n\norm{X_i}_1
  \]

  が成り立つ．
\end{lemma}

\begin{proof}
  許容分解 $X=\sum_jX^{(j)}$ を取る．$i\neq j$ なら，第 $i$ qubit だけを残す部分跡は
  第 $j$ qubit に関する部分跡を含むため，$\Tr_jX^{(j)}=0$ より寄与は $0$ である．したがって

  \[
    X_i=\Tr_{\{1,\ldots,n\}\setminus\{i\}}X^{(i)}.
  \]

  部分跡による跡ノルムの縮小（補題~\ref{lem:q-positive-contraction}）から
  $\norm{X_i}_1\leq\norm{X^{(i)}}_1$ である．$i$ について足し，許容分解の最小値を取ればよい．
\end{proof}

\begin{proposition}{計算基底上では Hamming 距離になる}{q-basis-hamming}
  任意の $x,y\in\{0,1\}^n$ に対し

  \[
    \Wass_1^{\rm q}(\proj{x},\proj{y})=d_H(x,y)
  \]

  が成り立つ．
\end{proposition}

\begin{proof}
  $x$ と $y$ が異なる座標を一つずつ変更する経路
  $x=x^{(0)},x^{(1)},\ldots,x^{(k)}=y$ を取る．ここで $k=d_H(x,y)$ である．
  隣り合う二点は一座標だけ異なるので，対応する純粋状態は隣接し，
  命題~\ref{prop:q-neighbor-form}より距離は $1$ 以下である．三角不等式から
  $\Wass_1^{\rm q}(\proj{x},\proj{y})\leq k$ を得る．

  逆に $X=\proj{x}-\proj{y}$ の第 $i$ 縮約は，$x_i=y_i$ なら $0$，
  $x_i\neq y_i$ なら跡ノルムが $2$ である．補題~\ref{lem:q-marginal-lower}より
  $\norm{X}_{W_1}\geq k$ だから，等号が成り立つ．
\end{proof}

\section{対角状態への還元}

\begin{theorem}{古典 Hamming--Wasserstein 距離の回復}{q-classical-recovery}
  $\{0,1\}^n$ 上の確率分布 $p,q$ と，対応する対角状態

  \[
    \rho_p=\sum_xp(x)\proj{x},
    \qquad
    \rho_q=\sum_yq(y)\proj{y}
  \]

  に対し

  \[
    \boxed{
      \Wass_1^{\rm q}(\rho_p,\rho_q)=\Wass_1^H(p,q)
    }
  \]

  が成り立つ．
\end{theorem}

\begin{proof}
  \textbf{量子距離が古典距離以下であること．}
  $\gamma\in\Pi(p,q)$ を任意の古典 coupling とすると

  \[
    \rho_p-\rho_q
    =\sum_{x,y}\gamma(x,y)(\proj{x}-\proj{y}).
  \]

  量子 $W_1$ ノルムの三角不等式と命題~\ref{prop:q-basis-hamming}より

  \[
    \norm{\rho_p-\rho_q}_{W_1}
    \leq\sum_{x,y}\gamma(x,y)d_H(x,y).
  \]

  $\gamma$ について最小値を取れば
  $\Wass_1^{\rm q}(\rho_p,\rho_q)\leq\Wass_1^H(p,q)$ を得る．

  \textbf{古典距離が量子距離以下であること．}
  $\rho_p-\rho_q=\sum_iX^{(i)}$ を量子 $W_1$ の許容分解とする．
  計算基底以外の成分を消す写像

  \[
    \Delta(A):=\sum_x\bra{x}A\ket{x}\proj{x}
  \]

  を各 $X^{(i)}$ に作用させる．$\Delta$ は半正定値性と跡を保つので，
  補題~\ref{lem:q-positive-contraction}より
  $\norm{\Delta(X^{(i)})}_1\leq\norm{X^{(i)}}_1$ である．また部分跡と可換だから
  $\Tr_i\Delta(X^{(i)})=0$ である．

  $\Delta(X^{(i)})$ の対角成分を符号付きベクトル $r^{(i)}$ とする．その正負部分の質量はともに
  $c_i:=\norm{r^{(i)}}_1/2$ である．$c_i>0$ のとき，正負部分を $c_i$ で割った確率分布は，
  第 $i$ 座標を除いた周辺分布が等しい．補題~\ref{lem:q-classical-neighbor}と
  Kantorovich--Rubinstein ノルムの斉次性より

  \[
    \norm{r^{(i)}}_{\rm KR}\leq c_i
    =\frac12\norm{\Delta(X^{(i)})}_1.
  \]

  さらに $p-q=\sum_ir^{(i)}$ だから，三角不等式を用いて

  \[
    \Wass_1^H(p,q)
    =\norm{p-q}_{\rm KR}
    \leq\sum_i\norm{r^{(i)}}_{\rm KR}
    \leq\frac12\sum_i\norm{X^{(i)}}_1.
  \]

  すべての許容分解について最小値を取れば逆向きの不等式を得る．
\end{proof}

以上で，量子 $W_1$ は量子状態全体で距離となり，古典部分では既存の Hamming--Wasserstein
距離を正確に再現することが分かった．本資料はここで終える．次の発展としては，目的に応じて
量子 Sinkhorn，coupling に基づく量子 $W_2$，または量子 Markov 半群の勾配流幾何へ進むのが自然である．
